\chapter{Induction}\label{chap:Induction}

\begin{fieldaxioms}
Addition and multiplication of real numbers are commutative and associative.
\end{fieldaxioms}

\begin{intuitivedef}\label{def:natural}
The natural numbers are the numbers that one gets from $1$ by adding $1$
again and again.
\end{intuitivedef}

\begin{theorem}[Principle of Mathematical Induction]\label{thm:PMI}
If $1$ is in a set $S$ of natural numbers, and $k+1$ is in $S$ whenever $k$
is, then $S$ holds every natural number.
\end{theorem}

\begin{skeleton}
Show the base case, then the inductive step, as in Theorem~\ref{thm:PMI}.
\end{skeleton}

\begin{problem}
Prove that $1 + 3 + \cdots + (2n-1) = n^2$ for each natural number $n$.
\end{problem}

\begin{question}
What goes wrong in a proof by induction with no base case?
\end{question}
